% Imporditud: tools/import_eksperimendid.py
% Allikas: Eesti füüsikaolümpiaadi eksperimendi ülesannete
%   kogu (Rael Kalda, 2023), ülesanne Ü136, lahendus L136.
\setAuthor{EFO žürii}
\setRound{piirkonnavoor}
\setYear{2009}
\setNumber{G E2}
\setDifficulty{4}
\setTopic{elektriahelad}
\setVahendid{Voltmeeter, takisti takistusega R = 7,5 k$\Omega$, värske tsink-süsinik patarei (sisetakistus alla 30 $\Omega$), ühendusjuhtmed}
\setPunktid{10}
\setAllikas{Eksperimendi kogu Ü136, lahendus lk 104}
\setLahendusseis{toores_ilma_skeemita}

\prob{Galvanomeeter}
Leida mitmele amprile vastab voltmeetri skaala maksimaalne väärtus, kui voltmeetrit kasutatakse ampermeetrina. Mõõtevea hindamist ei nõuta.

\hint

\solu
1) Mõõta vooluallika pinge. 2) Ühendada voltmeeter jadamisi takistiga. 3) Arvutada pingelang tuntud takistil mis on vooluallika pinge miinus voltmeetri näit jadaühenduses takistiga. 4) Ohmi seadusest leida voolutugevus tuntud takistil. 5) Teisendada voltmeetri näit ampritesse.

\textit{Žürii hindamisskeem on kogumikus lk 104; siia ei ole seda ümber kirjutatud.}
\probend
